\chapter{Công việc và trưởng thành}
\markboth{Công việc và trưởng thành}{Work and Maturity}

\section{Người mới đi làm nghĩ chỉ cần giỏi chuyên môn là đủ.}
\begin{truyen}{Người mới đi làm nghĩ chỉ cần giỏi chuyên môn là đủ.}{Skill is not the whole career.}
\chuhoa{C}{ậu làm việc tốt nhưng khó hợp tác, hay phản ứng và chê người khác chậm. Sau vài lần hụt cơ hội, cậu mới hiểu nghề nghiệp không chỉ cần năng lực làm việc mà còn cần năng lực làm người.}
\textit{(He worked well but was hard to cooperate with, reactive, and quick to dismiss others. After missing several opportunities, he learned that a career requires not only work skills but also human skills.)}
\end{truyen}
\begin{baihoc}
Giỏi việc giúp ta được chú ý; biết cư xử giúp ta được giữ lại lâu dài.
\textit{(Being skilled gets you noticed; knowing how to behave helps you stay.)}
\end{baihoc}
\begin{ghinhoanh}
\textbf{Short English to remember:}

Skill is not the whole career.

Competence needs character.
\end{ghinhoanh}
\begin{tuhocnhanh}
1. Em đang đầu tư nhiều hơn cho chuyên môn hay cho cách làm việc với người khác? \textit{(Are you investing more in technical skill or in how you work with people?)}

2. Điểm nào trong thái độ làm việc của em cần sửa sớm? \textit{(Which part of your work attitude should you fix early?)}
\end{tuhocnhanh}
\ngancach

\section{Một người nghĩ chăm chỉ là sẽ tự được nhìn thấy.}
\begin{truyen}{Một người nghĩ chăm chỉ là sẽ tự được nhìn thấy.}{Effort also needs expression.}
\chuhoa{A}{nh làm nhiều nhưng không biết báo cáo, không biết trình bày kết quả, không biết xin cơ hội đúng lúc. Khi người khác đi lên nhanh hơn, anh mới hiểu nỗ lực cần được thể hiện rõ ràng chứ không phải im lặng chịu thiệt mãi.}
\textit{(He worked hard but did not know how to report, present results, or ask for opportunity at the right time. When others advanced faster, he learned that effort must also be communicated clearly, not hidden in silence forever.)}
\end{truyen}
\begin{baihoc}
Làm tốt là nền, nhưng nói cho đúng và đúng lúc cũng là một phần của trưởng thành nghề nghiệp.
\textit{(Doing good work is the foundation, but communicating it well and at the right time is also part of professional maturity.)}
\end{baihoc}
\begin{ghinhoanh}
\textbf{Short English to remember:}

Effort also needs expression.

Silent work can stay unseen.
\end{ghinhoanh}
\begin{tuhocnhanh}
1. Em có đang làm tốt nhưng trình bày kém không? \textit{(Are you doing well but communicating it poorly?)}

2. Em cần học kỹ năng nào để nỗ lực của mình được nhìn đúng hơn? \textit{(What skill do you need so your effort is seen more accurately?)}
\end{tuhocnhanh}
\ngancach

\section{Một người nhảy việc liên tục vì chỗ nào cũng thấy không vừa ý.}
\begin{truyen}{Một người nhảy việc liên tục vì chỗ nào cũng thấy không vừa ý.}{No place will grow you if you flee every discomfort.}
\chuhoa{C}{hỗ nào có áp lực là anh muốn rời. Chỗ nào có người khó tính là anh muốn đổi. Sau nhiều lần bắt đầu lại, anh mới hiểu không phải nơi nào cũng sai; có khi mình đang trốn chính phần cần trưởng thành nhất.}
\textit{(Whenever work became hard or someone demanding appeared, he wanted to leave. After many restarts, he realized not every place was wrong; sometimes he was running from the part of himself that needed growth most.)}
\end{truyen}
\begin{baihoc}
Không phải mọi khó chịu đều là dấu hiệu phải rời đi; có cái là dấu hiệu mình cần lớn lên.
\textit{(Not every discomfort is a sign to leave; some are signs that we need to grow.)}
\end{baihoc}
\begin{ghinhoanh}
\textbf{Short English to remember:}

Do not flee every discomfort.

Some pressure is growth in disguise.
\end{ghinhoanh}
\begin{tuhocnhanh}
1. Em có đang rời đi quá nhanh khỏi một nơi có thể dạy mình điều cần học không? \textit{(Are you leaving too quickly from a place that could teach you something important?)}

2. Khó chịu nào của hiện tại thật ra là bài học? \textit{(Which present discomfort might actually be a lesson?)}
\end{tuhocnhanh}
\ngancach

\section{Một người chỉ hiểu giá trị của kỷ luật sau nhiều năm sống theo cảm hứng.}
\begin{truyen}{Một người chỉ hiểu giá trị của kỷ luật sau nhiều năm sống theo cảm hứng.}{Discipline saves scattered years.}
\chuhoa{N}{gày hứng thì làm nhiều, ngày chán thì bỏ ngang. Những năm tháng như thế trôi qua mà thành quả ít ỏi. Khi học được nhịp đều, anh mới hiểu kỷ luật không giết tự do; nó cứu tự do khỏi tan vào lười biếng.}
\textit{(On inspired days he worked a lot, on dull days he quit halfway. Years passed with little to show. When he learned consistency, he understood that discipline does not kill freedom; it saves freedom from dissolving into laziness.)}
\end{truyen}
\begin{baihoc}
Nhiều người hiểu muộn rằng cảm hứng đẹp, nhưng kỷ luật mới là thứ đưa đời mình đi xa.
\textit{(Many people learn too late that inspiration is beautiful, but discipline is what carries life further.)}
\end{baihoc}
\begin{ghinhoanh}
\textbf{Short English to remember:}

Discipline saves scattered years.

Consistency beats mood.
\end{ghinhoanh}
\begin{tuhocnhanh}
1. Việc nào trong đời em đang bị tâm trạng chi phối quá nhiều? \textit{(What part of your life is being ruled too much by mood?)}

2. Một nếp kỷ luật nhỏ nào em có thể dựng ngay? \textit{(What small discipline can you build right now?)}
\end{tuhocnhanh}
\ngancach

\section{Người trưởng thành hiểu muộn rằng không ai nợ mình cơ hội mãi.}
\begin{truyen}{Người trưởng thành hiểu muộn rằng không ai nợ mình cơ hội mãi.}{No one owes you endless chances.}
\chuhoa{C}{ó lần đầu người ta bỏ qua, lần hai người ta nhắc nhẹ, lần ba người ta bắt đầu rút niềm tin lại. Anh từng nghĩ người khác sẽ mãi cho mình cơ hội, cho đến khi cánh cửa khép hẳn vì những lần xem nhẹ liên tiếp.}
\textit{(The first time people overlook it, the second time they remind gently, the third time trust starts to fade. He once thought others would always give him another chance, until a door fully closed after repeated carelessness.)}
\end{truyen}
\begin{baihoc}
Được tin lại là một ân huệ, không phải điều mặc nhiên phải có mãi.
\textit{(Being trusted again is a grace, not something guaranteed forever.)}
\end{baihoc}
\begin{ghinhoanh}
\textbf{Short English to remember:}

No one owes you endless chances.

Treat trust carefully.
\end{ghinhoanh}
\begin{tuhocnhanh}
1. Em có đang dùng sự tin tưởng của ai quá dễ dãi không? \textit{(Are you treating someone's trust too carelessly?)}

2. Cơ hội nào em cần trân trọng hơn ngay lúc này? \textit{(What opportunity do you need to value more right now?)}
\end{tuhocnhanh}
\ngancach